\subsection{Threats to validity}

\begin{itemize}
  \item \textbf{Implementation scope.} The reference implementation is a
  standalone Python package; it exercises the gate logic, not a full agent
  harness. Integration with a live long-horizon agent is future work.
  \item \textbf{Formal statements vs.\ certified systems.} The Lean files
  state the theorems and prove their structural forms; they do not certify
  the Python implementation. The two artifacts are related by statement,
  not by a verified compiler.
  \item \textbf{Adversarial coverage.} The evaluation battery covers the
  ten transition fixtures plus resume semantics. It is not a fuzzing
  campaign, a penetration test, or a formal verification of the gate.
  \item \textbf{No live middle-layer measurement.} We do not claim
  middle-layer intelligence verification; no live activation capture was
  performed for this paper.
  \item \textbf{No production certification.} Nothing here certifies a
  product, a regulatory posture, a court-admissibility property, or a
  hardware root of trust.
\end{itemize}

\subsection{Claim ceiling}

\begin{center}
\fbox{\parbox{0.85\textwidth}{\centering
OPEN RESEARCH ARTIFACT\\
LOCAL AND CI-REPRODUCIBLE SOFTWARE EVIDENCE\\
NOT PRODUCT, REGULATORY, OR SAFETY CERTIFICATION}}
\end{center}

Results apply only to the released configuration
(\texttt{ProofOS/research/no-witness-no-verdict}, tagged
\texttt{v0.1.0-research}). The released artifact identifies source commit,
tree digest, artifact SHA-256, Python and Lean versions, the exact test
command, the exact test summary, and known limitations.
